\clearpage
\subsection{Europa}

\subsubsection{Unión Europea (UE)}

La \acrfull{ue} es el ejemplo por excelencia de un proceso de integración que culminó en una unión económica y monetaria. Sus antecedentes se remontan a la creación de la \gls{ceca} en 1951 por seis países europeos: Francia, Italia, Bélgica, Luxemburgo, Países Bajos y la República Federal de Alemania. En 1957, esos mismos países formaron la \gls{cee} a través de la firma del Tratado de Roma. Entre los principales elementos de la creación de la \gls{cee} destacan la constitución de una unión aduanera que eliminó los aranceles y las cuotas de importación entre los Estados miembros, la adopción de un arancel externo común para las importaciones procedentes de terceros países, así como el desarrollo de políticas comunes en materia agrícola, de transporte y de competencia \parencite[p.~180]{Leuffen2022}. En paralelo, ese mismo año también se constituyó la \gls{euratom} mediante un tratado independiente para promover el crecimiento de la industria nuclear en Europa. En 1965, las instituciones de las tres comunidades europeas se fusionaron para formar una sola Comisión Europea mediante la firma del Tratado de Bruselas \parencite[p.~563]{EuropaDirectory2022}. Los siguientes 15 años fueron de una intensa integración dentro de la \gls{cee}, acompañada de altas tasas de crecimiento económico y las primeras políticas comunes, como la Política Agrícola Común de 1962 \parencite[p.~48]{Kollarik2025}.

%La consolidación del mercado común también fue impulsada por la jurisprudencia del Tribunal de Justicia de las Comunidades Europeas. Mediante los casos \textit{Van Gend en Loos} de 1963 y \textit{Costa contra ENEL} de 1964, el Tribunal estableció, respectivamente, los principios de efecto directo y primacía del derecho comunitario, lo que permitió a empresas y particulares invocar sus derechos ante los tribunales nacionales e impidió que los Estados miembros contravinieran unilateralmente las normas comunes del mercado \parencite[pp.~2413--2415]{Weiler1991}. https://scispace.com/pdf/the-transformation-of-europe-1cnmksqq67.pdf

Debido a las crisis económicas durante la década de 1970, el proceso de integración europea se ralentizó durante los siguientes 15 años. Pese a ello, la \gls{cee} tuvo su primera ampliación en 1973 con la adhesión del Reino Unido, Dinamarca e Irlanda; después se unió Grecia en 1981 y España y Portugal en 1986 \parencite[pp.~48--49]{Kollarik2025}. A partir de la reunificación de Alemania en 1990, la antigua República Democrática Alemana pasó a formar parte de la Comunidad de forma inmediata. Cinco años después, Austria, Finlandia y Suecia se convirtieron en miembros \parencite[p.~564]{EuropaDirectory2022}.

En 1985, los Estados miembros de la \gls{cee} reconocieron que el nacionalismo económico y la falta de cooperación política afectaban negativamente la competitividad europea en el marco mundial. Por lo tanto, iniciaron un proceso de implementación de un mercado común con la firma del Acta Única Europea en 1986. El establecimiento del mercado común también requirió que la comunidad comenzara a planificar el avance hacia una unión monetaria \parencite[p.~49]{Kollarik2025}.

A partir de 1993, la \gls{cee} pasó a denominarse formalmente \gls{ce} con la entrada en vigor del Tratado de Maastricht. Este tratado creó la \acrfull{ue} con el objetivo de fortalecer la cooperación económica y monetaria, introducir la cooperación en materia de justicia, establecer una política exterior y de seguridad común e introducir la ciudadanía europea \parencite[p.~563]{EuropaDirectory2022}. El tratado también significó la transformación del mercado común en un mercado único y la implementación de la \gls{uem} en tres etapas, que culminaría con la introducción de la moneda única (el euro) en 11 países miembros en 1999 \parencite[p.~598]{EuropaDirectory2022}. 

De acuerdo con \cite{Kollarik2025}, el establecimiento de la \gls{uem} convirtió a la Unión Europea en la integración regional más avanzada de la economía mundial y en el mayor bloque comercial en el mundo. Sin embargo, la unión económica y monetaria creó una estructura de integración diferenciada, ya que no todos los países miembros adoptaron la moneda única inmediatamente. Además, la política monetaria supranacional no fue complementada con una política fiscal ni una regulación financiera común, ya que las políticas presupuestarias y la regulación de instituciones financieras siguieron estando bajo la autoridad exclusiva de los gobiernos nacionales. Estas ausencias institucionales se hicieron evidentes después del estallido de la crisis financiera de 2008, lo que conllevó a una reestructuración del sistema de la \gls{uem}, principalmente en el ámbito fiscal y financiero \parencite[pp.~55--56]{Kollarik2025}.

El euro entró en vigor el 1 de enero de 1999 en 11 Estados miembros\footnote{Austria, Bélgica, Finlandia, Francia, Alemania, Irlanda, Italia, Luxemburgo, Países Bajos, Portugal y España} de la Unión Europea, conformando lo que se conoce como la eurozona. Posteriormente, se integraron otros 10 miembros\footnote{Grecia, Eslovenia, Chipre, Malta, Eslovaquia, Estonia, Letonia, Lituania, Croacia y Bulgaria.} a la Eurozona; por su parte, el resto de Estados miembros,\footnote{Hungría, Polonia, República Checa, Rumania y Suecia.} salvo Dinamarca, que cuenta con una cláusula de exclusión, están obligados a adoptar también el euro \parencite[p.~51]{Kollarik2025}. Otros Estados no miembros, como Montenegro y Kosovo, junto con algunos microestados europeos, han adoptado el euro como moneda, sin participar en el Sistema Europeo de Bancos Centrales \parencite[p.~2]{Leuffen2022}.

En 2004, ocho países de Europa central y oriental,\footnote{República Checa, Eslovaquia, Hungría, Eslovenia, Polonia, Estonia, Letonia y Lituania.} junto con Malta y Chipre, se incorporaron a la Unión Europea; tres años después también lo harían Bulgaria y Rumania. Finalmente, el último país en integrarse a la \gls{ue} fue Croacia en 2013, con lo que la Unión Europea alcanzó un total de 28 Estados miembros. Otros países han solicitado su adhesión o se encuentran en negociaciones para ello.\footnote{Turquía, Macedonia del Norte, Albania, Montenegro, Serbia, Bosnia y Herzegovina, Ucrania, Georgia y Moldavia} La \gls{ue} experimentó la primera salida de un Estado miembro con el caso del Reino Unido, que mediante un referéndum nacional celebrado en 2016, conocido comúnmente como Brexit, decidió abandonar la \acrlong{ue}, decisión que se concretó en 2020 \parencite[pp.~564--565]{EuropaDirectory2022}. 

Existen otros países no miembros que participan en diversos grados en los marcos comerciales de la \gls{ue}. Tal es el caso del \gls{eee}, que comprende el mercado interior común y políticas complementarias que se extienden a tres de los cuatro países que conforman la \gls{efta}.\footnote{Islandia, Noruega y Liechtenstein. Suiza, aunque forma parte de la \gls{efta}, no se integró al \gls{eee}; no obstante, ha suscrito varios tratados bilaterales con la \gls{ue} y ha incorporado normas comunitarias individuales en su legislación nacional.} Por su parte, la unión aduanera también incluye a Turquía. Además, algunos países de Europa del Este, Medio Oriente y África del Norte que participan en la Política Europea de Vecindad adoptan algunos elementos de las normas comunitarias, aunque no tienen oportunidad de unirse a la \gls{ue} a corto plazo \parencite[p.~4]{Leuffen2022}. En contraste, la unión monetaria está diferenciada al interior de la \gls{ue}, debido a que hay algunos miembros que no están dispuestos o no pueden unirse a la zona euro \parencite[p.~381]{Leuffen2022}. 

El comercio intracomunitario de la \gls{ue} se rige por la política de la unión aduanera que comenzó a funcionar en 1968. Como resultado, el comercio entre los Estados miembros quedó libre de aranceles y de otras restricciones no arancelarias. A su vez, se estableció una política arancelaria común hacia terceros países. El alto nivel de integración y la proximidad geográfica han conducido a un comercio intrarregional muy intenso. Una muestra de ello es que en la mayoría de los países miembros de la \gls{ue}, las exportaciones hacia otros países de la unión representan entre el 50\% y el 70\% de sus exportaciones totales \parencite[pp.~51--52]{Kollarik2025}.

Respecto a la política comercial exterior, la \gls{ue} ha establecido diversos acuerdos comerciales y económicos con casi todas las regiones del mundo. Entre estos instrumentos, los \gls{aa} constituyen uno de los principales mecanismos mediante los cuales la \gls{ue} articula sus relaciones con el exterior. Los \glspl{aa} han servido como instrumentos de preadhesión de países candidatos; por ejemplo, los acuerdos de asociación y estabilización con los países balcánicos. En segundo lugar, también han servido como una alternativa a la membresía plena; este es el caso del Acuerdo sobre el \gls{eee} y los \glspl{aa} con países de Europa del Este como Ucrania, Moldavia y Georgia. En último lugar, los \glspl{aa} también han servido como marco para relaciones privilegiadas con otras regiones, como los Acuerdos Euromediterráneos con países del norte de África y los Acuerdos de Asociación con América Central y América Latina. Además de estos mecanismos, la \gls{ue} cuenta con numerosos acuerdos de libre comercio con países de renta media y alta de todo el mundo \parencite[p.~57]{Kollarik2025}. 

Uno de los elementos más importantes de la integración de la \gls{ue} fue que el mercado único no solo implicó la libre circulación de mercancías y capitales, sino también de trabajadores y personas mediante la reducción de las barreras físicas, técnicas y fiscales \parencite[p.~53]{Kollarik2025}. El área de libre circulación europea comenzó con la firma del Acuerdo de Schengen en 1985, fuera del marco comunitario, entre cinco países de la \gls{ce}.\footnote{Francia, Alemania, Bélgica, Luxemburgo y los Países Bajos} En la siguiente década, estas políticas fueron integradas a la \gls{ue} y se extendieron a todos sus miembros, con la excepción de Irlanda, el Reino Unido, Chipre, Bulgaria y Rumania. Por el contrario, hay cuatro países no miembros de la \gls{ue} que pertenecen al espacio Schengen: Islandia, Noruega, Suiza y Liechtenstein, junto con otros cuatro micro-Estados europeos \parencite[p.~3]{Leuffen2022}.

\subsubsection{Asociación Europea de Libre Comercio (EFTA)}

La creación de la \acrfull{cee} no tuvo un apoyo generalizado en Europa. En respuesta a su creación, Austria, Dinamarca, Noruega, Portugal, Suecia, Suiza y el Reino Unido fundaron la \acrfull{efta} en 1960. Posteriormente, se unieron a esta asociación Finlandia en 1961,\footnote{En 1961 ingresó como miembro asociado y hasta 1986 se convirtió en miembro pleno.} Islandia en 1970 y Liechtenstein en 1991 \parencite[p.~181]{Leuffen2022}.
En 1966, los Estados miembros de la \gls{efta} alcanzaron el libre comercio total de productos industriales y, durante la década de 1970, suscribieron acuerdos bilaterales de libre comercio con la \gls{cee} \parencite{EFTA}. En 1973, el Reino Unido y Dinamarca decidieron abandonar la asociación para ingresar a la \gls{cee}, y Portugal siguió el mismo camino en 1986. Finalmente, con la entrada de Austria, Finlandia y Suecia a la \gls{ue} y su salida de la \gls{efta}, la importancia de esta asociación disminuyó significativamente \parencite[p.~69]{Leuffen2022}. Como resultado, la \gls{efta} quedó integrada únicamente por Noruega, Islandia, Liechtenstein y Suiza, países que permanecieron fuera de la \gls{ue}. Desde la década de 1990, la \gls{efta} desarrolló una política de negociación conjunta de acuerdos de libre comercio con terceros países. Los Estados miembros de la \gls{efta} también desarrollaron estrechas relaciones comerciales con la UE, como se refleja en el acuerdo de creación del \acrfull{eee} en 1994 y los acuerdos bilaterales entre Suiza y la \gls{ue} de 1999 \parencite{EFTA_Convention_overview}. 


El Convenio de la \gls{efta}, revisado en 2001, regula las relaciones económicas entre los cuatro Estados miembros de la asociación, abarcando el comercio de bienes, servicios e inversiones, el transporte aéreo y terrestre, la agricultura, los derechos de propiedad intelectual, la contratación pública, la libre circulación de personas, la seguridad social, las barreras técnicas al comercio, la competencia, las empresas públicas y las ayudas estatales \parencite{EFTA_Convention}. 

%En la década de 1990, los Estados de la \gls{efta} suscribieron trece \gls{tlc}, cuatro de los cuales aún conintúan vigentes (Turquía, Marruecos, Israel y la Autoridad Palestina); mientras que los restantes, suscritos con países de Europa Central y Oriental, dejaron de estar en vigor tras la adhesion de estos países a la \gls{ue}. Posteriormente, se firman varios acuerdos con otros países e instancias, entre los que destaca el acuerdo de libre comercio con \gls{MERCOSUR} en 2025 y \parencite{EFTA}.

La \gls{efta} se constituyó como una zona de libre comercio, con poca burocracia y sin expectativas de contribución financiera de sus miembros \parencite[p.~5]{Randwyck2011} y, a diferencia de la \gls{cee}, la \gls{efta} no estableció un arancel aduanero externo común \parencite[p.~181]{Leuffen2022}. El comercio de mercancías entre los miembros de la \gls{efta} se caracteriza por el intercambio de productos de alta intensidad de capital, como productos farmacéuticos, barcos y maquinaria eléctrica, así como de recursos naturales, principalmente minería y combustibles \parencite{EFTA_Convention}.

Aunque en términos políticos, jurídicos e institucionales la \gls{efta} se mantuvo independiente de la \gls{cee} y, después, de la \gls{ue}, la integración comercial de los miembros de la \gls{efta} con la comunidad europea ha sido profunda. A partir de 1984, se eliminaron las últimas barreras arancelarias entre ambos grupos y se estableció un tratado de libre comercio total para los productos industriales entre la \gls{cee} y los países de la \gls{efta}. La entrada en vigor del \gls{eee} en 1994 permitió a Noruega, Islandia y Liechtenstein unirse a un mercado único europeo de bienes, servicios, capital y trabajo, sin necesidad de adoptar las instituciones políticas de la \gls{ue} ni sus políticas comunes en agricultura, pesca, justicia, moneda única o política exterior común \parencite[p.~5]{Leuffen2022}. En cambio, Suiza optó por establecer acuerdos bilaterales de libre comercio y numerosos acuerdos sectoriales con la \gls{ue} \parencite[p.~604]{EuropaDirectory2022}. 

\subsubsection{Análisis de comunidades}

A pesar del alto grado de integración comercial y económica alcanzado en Europa por medio de la \gls{ue}, la \gls{efta} y el \gls{eee}, los resultados de la detección de comunidades muestran una estructura fragmentada en varios bloques subregionales. Las comunidades identificadas en Europa pueden agruparse en cuatro tipos:
\begin{itemize}
	\item Comunidades subregionales estables, integradas de manera continua por las economías de varios países europeos. En este grupo se encuentran las comunidades de Europa central, los países nórdicos y bálticos, la región Benelux\footnote{Bélgica, Países Bajos y Luxemburgo.}, así como el bloque formado por el Reino Unido e Irlanda. 
	\item Comunidades multinacionales que agrupan economías europeas y países africanos, especialmente antiguas colonias. Este patrón se observa en la comunidad encabezada por Francia y en la integrada por España y Portugal.
	\item Comunidades inestables, caracterizadas por cambios frecuentes en su composición y en la pertenencia de sus miembros. Este es el caso de las economías de los Balcanes (Grecia, Bulgaria y Rumania), de Europa del Este (Rusia, Ucrania y Bielorrusia), así como de Turquía y Chipre.
	\item Comunidades constituidas por un solo país. Los casos más representativos son Suiza e Italia.
\end{itemize}
En la figura \ref{Mapa_Comunidades_Europa} se muestran las subgráficas inducidas por los vértices de cada comunidad en Europa en tres años específicos: 1995, 2008 y 2022. Para facilitar su visualización, solo se representan los flujos comerciales más importantes dentro de cada bloque. La figura permite identificar con claridad varias comunidades subregionales estables, cuya composición se mantiene a lo largo del tiempo e incluso se amplía. Este es el caso del bloque productivo de Europa central y del conformado por los países nórdicos y bálticos. En contraste, los países balcánicos y las antiguas repúblicas soviéticas europeas presentan cambios en su pertenencia comunitaria durante los años mostrados.

Entre las comunidades subregionales estables destaca el caso de Europa central. Entre 1995 y 1998, las economías completas de Austria, República Checa, Hungría, Eslovaquia, Eslovenia y Croacia conformaron una comunidad, mientras que Alemania constituyó otro por separado. A partir de 2000, ambos bloques se fusionaron en una sola comunidad encabezada por Alemania. Este proceso resulta relevante porque la integración estable de varias economías centroeuropeas en un mismo bloque productivo antecedió a su incorporación a la \gls{ue} en 2004, como ocurrió con República Checa, Eslovaquia, Eslovenia y Hungría. Polonia, que también ingresó a la \gls{ue} ese año, se incorporó a la comunidad en 2010. Croacia, por su parte, se separó del bloque y permaneció como una comunidad individual entre 1999 y 2012, con excepción de 2009; posteriormente, volvió a integrarse al bloque encabezado por Alemania en 2013, mismo año en el que se formalizó su ingreso a la \gls{ue}. Rumania, que desde 2007 era parte de la \gls{ue}, pasó a formar parte de este bloque después de 2015. En la figura \ref{ComunidadesPorPais_DEU} se puede observar a detalle el proceso de consolidación de este bloque productivo. 

\begin{figure}[H]
	\centering
	\includegraphics[width=1\textwidth]{images/ComunidadesPorPais_DEU.png}
	\caption{Integración de las economías de Europa central en un bloque productivo.}
	\label{ComunidadesPorPais_DEU}
\end{figure}

En conjunto, los resultados muestran la consolidación progresiva de un bloque productivo centroeuropeo, inicialmente integrado por un núcleo constante de economías y fortalecido posteriormente mediante la incorporación de nuevos países. La estabilidad en el tiempo de este bloque evidencia que los vínculos productivos detectados por el análisis de comunidades son estructurales, y no responden sólo a fluctuaciones coyunturales. Destaca también que en varios casos, la pertenencia a la comunidad ocurre antes o coincide con su adhesión a la UE, lo que sugiere que la convergencia económica no fue únicamente una consecuencia de la membresía, sino que ya existían vínculos productivos significativos. 

La permanencia de Alemania como líder dentro de la comunidad de Europa central coincide con la literatura que identifica a Alemania como el centro industrial de la denominada \textit{Factory Europe}, desde el cual se articula la producción y se redirigen las exportaciones de las economías de Europa central y oriental hacia mercados fuera de Europa, principalmente hacia Estados Unidos, Rusia y China \parencite{KordalskaOlczyk2019}. Esta interpretación también se refleja en la figura \ref{Mapa_Comunidades_Europa}, donde, en 2008 y 2022, los flujos comerciales más intensos dentro de la comunidad conectan a la economía alemana principalmente con las economías checa, austriaca y polaca y, en menor medida, con las de Hungría y Eslovaquia.

En la figura \ref{Subgraph_deu}, que presenta la subgráfica inducida por los países de Europa central en 2022, tanto la posición de Alemania, obtenida mediante el algoritmo de Fruchterman--Reingold (véase el anexo \ref{anexo5}), como el tamaño de sus vértices refuerzan el argumento sobre el papel central que desempeña la economía alemana en la región. Asimismo, se observa cierto equilibrio entre los flujos comerciales intrarregionales y los internos, particularmente en las relaciones entre los sectores económicos de Alemania, Austria, Polonia, República Checa, Hungría y Rumanía. Estos flujos no solo se dirigen hacia Alemania, sino que también conectan a dichas economías entre sí. En contraste, Croacia parece presentar un menor grado de cohesión con el bloque subregional.

\begin{figure}[H]
	\centering
	\includegraphics[width=1\textwidth]{subgraphs/deu_subrgraph.png}
	\caption{Subgráfica inducida por los vértices de los países de Europa central (2022).$^{*}$}
	\label{Subgraph_deu}
	
	\caption*{\footnotesize $^{*}$ El tamaño de cada vértice representa su fuerza de salida; su posición se obtiene mediante el algoritmo de Fruchterman--Reingold (500 iteraciones); y su color identifica el país al que pertenece, no su comunidad. Solo se presentan las aristas cuyo peso supera el percentil 95 de la distribución regional. Las aristas grises representan transacciones domésticas y las de color coral, flujos comerciales internacionales.}
\end{figure}

Otro ejemplo de una comunidad estable identificada en Europa es el caso de los países nórdicos y bálticos. Esta comunidad inicia agrupando las economías de los países nórdicos (Dinamarca, Noruega, Finlandia y Suecia), así como los sectores de Estonia e Islandia. Todos estos países se mantienen dentro de la comunidad durante todo el periodo estudiado. La comunidad es identificada con la economía de Suecia debido a que es el país con el sector con el mayor flujo comercial de salida dentro de la comunidad. A partir de 2007,  la economía de Lituania comienza a aparecer enlazada a este bloque (con la excepción de los años 2008, 2012 y 2014); y Letonia sigue este camino desde 2015. Antes de eso, ambos países bálticos habían aparecido consistentemente dentro de la comunidad encabezada por la economía rusa, pese a haber ingresado a la Unión Europea en 2004. Finalmente, en 2022 el bloque productivo de los países nórdicos y bálticos parece fusionarse junto a la comunidad de Bélgica y Países Bajos. La evolución detallada de esta comunidad se muestra en la figura \ref{DiagramaSankey_Europa_Nordica}. 

\begin{figure}[H]
	\centering
	\includegraphics[width=1\textwidth]{sankey_diagrams/europe_nordic.png}
	\caption{Evolución de la pertenencia a comunidades de los sectores económicos de los países nórdicos y bálticos.}
	\label{DiagramaSankey_Europa_Nordica}
\end{figure}

Este bloque productivo presenta una elevada estabilidad temporal, pues conserva durante todo el periodo un núcleo constante integrado por las economías nórdicas. Como se observa en la figura \ref{Mapa_Comunidades_Europa}, sus principales flujos comerciales muestran una estructura relativamente equilibrada entre Dinamarca, Noruega y Suecia. A esta relación tripartita se añade la intensa vinculación de Finlandia con la economía sueca. La incorporación posterior de Lituania y Letonia sugiere, además, una ampliación gradual del espacio productivo hacia el Báltico y una reorientación de ambas economías desde el bloque encabezado por Rusia hacia sus socios del norte de Europa. El desfase entre su ingreso a la \gls{ue} en 2004 y su vinculación más estable con esta comunidad muestra que la integración institucional no se tradujo de manera inmediata en una reconfiguración de sus relaciones productivas. Finalmente, la unión observada en 2022 con la comunidad de Bélgica y Países Bajos podría indicar una articulación más amplia con el espacio productivo del norte de Europa, aunque, al tratarse de un cambio registrado únicamente al final del periodo, esta unión podría no ser significativa.


A nivel individual, algunos sectores vinculados con la energía, la minería y los metales de Noruega, Islandia y Lituania estuvieron temporalmente vinculados con comunidades distintas de la nórdico-báltica, como se observa en la figura \ref{DiagramaSankey_Europa_Nordica}. Esto sugiere que su articulación productiva respondió más a cadenas de valor específicas que a la proximidad regional. En Noruega, los sectores de extracción de petróleo crudo y gas natural (B06), servicios de apoyo a la minería (B09) y productos de refinación de petróleo (C19) estuvieron asociados con la comunidad encabezada por el Reino Unido entre 2005 y 2012, con excepción de 2009; además, el sector B06 mantuvo esta configuración entre 2014 y 2021. De forma similar, la permanencia de los sectores B09 y C19 de Lituania en el bloque liderado por Rusia entre 2015 y 2019 muestra que ciertos vínculos productivos heredados de la ex república soviética persistieron incluso después de que el conjunto de la economía comenzó a orientarse hacia el norte de Europa. Por último, la industria de metales preciosos básicos (C24B) de Islandia apareció integrada en el bloque productivo de los países centroeuropeos durante el periodo 2006--2022 (con la excepción de 2014 y 2016).
% Noruega - b06, b09, c19 - GB 2005-2012 (con excepcion de 2009 ) y otra vez en 2021
%  Noruega b06  - GBR 2014 a 2021
%  Islandia C24b - DEU 2006-2022 (excpcion 2014 y 2016)
%Lituania b09 y c19 - Rusia (2015 al 2019). En 2021 otra ves b09

Bélgica y los Países Bajos muestran una fuerte interdependencia productiva, ya que sus economías formaron una misma comunidad durante todo el periodo estudiado. Un patrón similar se observó entre Gran Bretaña e Irlanda. Además, entre 2015 y 2021, ambas comunidades se fusionaron en un solo bloque productivo, como se muestra en la figura \ref{ComunidadesPorPais_GBR}. Durante ciertos periodos, las economías de Luxemburgo y Malta también se integraron a estos bloques. Luxemburgo formó parte de la comunidad de Bélgica y los Países Bajos entre 1995 y 2005. Estos tres países conforman la unión de cooperación intergubernamental conocida como \textit{Benelux Union}. Posteriormente, entre 2006 y 2014, la economía de Luxemburgo parece haberse vinculado más estrechamente con el bloque productivo de Europa central. Por su parte, Malta cambió con frecuencia de comunidad hasta 2012. A partir de ese año, pasó a formar parte de manera estable de la comunidad de Europa noroccidental, encabezada por Gran Bretaña. La salida de Gran Bretaña de la \gls{ue} no parece haber afectado inmediatamente a este bloque productivo, aunque podría ser una de las razones por las que se dividió en dos en 2022 (véase la figura \ref{ComunidadesPorPais_GBR}).

\begin{figure}[H]
	\centering
	\includegraphics[width=1\textwidth]{images/ComunidadesPorPais_Benelux.png}
	\caption{Integración de las economías de Europa noroccidental en comunidades.}
	\label{ComunidadesPorPais_GBR}
\end{figure}

El análisis de comunidades mostró la existencia de bloques multinacionales integrados por economías europeas y africanas, principalmente de países mediterráneos y antiguas colonias. Uno de estos casos corresponde a España y Portugal, cuyas economías formaron una comunidad durante todo el periodo de estudio. Marruecos se incorporó en 2002. Por su parte, Angola y Santo Tomé y Príncipe, antiguas colonias portuguesas, estuvieron vinculadas a este bloque productivo durante varios periodos. Además, la extracción de petróleo y gas natural de Camerún (B06) también formó parte de esta comunidad entre 2011 y 2015, al igual que la silvicultura y la tala de la República Democrática del Congo (A02) entre 1995 y 1996 y entre 2002 y 2003. La figura \ref{ESP_timeline} muestra la cantidad de sectores económicos de cada país que formaron parte de la comunidad encabezada por España.


\begin{figure}[H]
	\centering
	\includegraphics[width=\linewidth]{images/communities_membership/ESP.png}
	\caption{Cantidad de sectores económicos (vértices) por año en la comunidad ESP (España).}
	\label{ESP_timeline}
\end{figure}

Algo similar ocurre en el caso francés. Durante la mayor parte del periodo de estudio, los sectores económicos de Francia y Túnez, antiguo protectorado francés, formaron parte del mismo bloque. En determinados años también se incorporaron otras antiguas colonias francesas: Marruecos entre 1997 y 2001, Camerún en 1997 y 2000, y Costa de Marfil en 2000 (véase la figura \ref{FRA_timeline}). Aunque estas economías africanas no permanecieron de manera continua en la comunidad, algunos de sus sectores sí se integraron al bloque encabezado por Francia durante periodos más prolongados. Entre ellos se encuentran los sectores mineros de Senegal (B05, B07 y C24B), que pertenecieron a esta comunidad entre 1996 y 2006; la extracción de metales preciosos de Camerún (C24B), integrada entre 1995 y 2011; y varios sectores industriales avanzados de Marruecos. En este último caso, la producción de computadoras y productos electrónicos (C26) formó parte del bloque entre 2002 y 2005; la fabricación de otros equipos de transporte (C31T33), a partir de 2016; y los servicios de tecnología e información (J62\_63), durante varios años desde 2009.

El hecho de que, a partir de 2000, las economías de estas antiguas colonias francesas dejaran de formar parte del bloque productivo encabezado por Francia podría indicar un debilitamiento de sus vínculos de dependencia y una reorientación de sus estructuras económicas hacia otros centros productivos. Desde 2001, Camerún y Costa de Marfil se integraron a una comunidad junto con Senegal y Nigeria. Un año después, la economía de Túnez pasó a formar parte del bloque productivo de España y Portugal (véase la figura \ref{ComunidadesPorPais}).

\begin{figure}[H]
	\centering
	\includegraphics[width=\linewidth]{images/communities_membership/FRA.png}
	\caption{Cantidad de sectores económicos (vértices) por año en la comunidad FRA (Francia).}
	\label{FRA_timeline}
\end{figure}

Los países de los Balcanes y las ex república soviéticas presentan la mayor inestabilidad en la región en cuanto a su pertenencia comunitaria. Durante buena parte del periodo de estudio, sus economías se distribuyeron entre tres comunidades encabezadas, respectivamente, por Rusia, Grecia y Turquía, como se muestra en las figuras \ref{Mapa_Comunidades_Europa} y \ref{ComunidadesPorPais_Rusia}.

A finales de la década de 1990, el bloque productivo encabezado por la Federación Rusa estaba integrado de manera relativamente estable por varias antiguas repúblicas soviéticas (Bielorrusia, Ucrania, Kazajistán, Lituania y Letonia), además de Bulgaria y Polonia. Como se observa con detalle en la figura \ref{RUS_timeline}, la composición de este bloque comenzó a cambiar durante la década siguiente. Bulgaria salió de la comunidad después de 2002 y se reincorporó en 2021. Por su parte, Polonia se integró al bloque productivo de Europa central a partir de 2006. Posteriormente, desde 2015, Lituania y Letonia pasaron a formar parte de la comunidad nórdico-báltica encabezada por Suecia. Estos cambios reflejan un debilitamiento de los vínculos productivos de Rusia con su antigua zona de influencia en el Báltico y Europa del Este. En contraste, pese a los conflictos políticos y militares entre Rusia y Ucrania posteriores a 2014, ambas economías mantuvieron una fuerte interdependencia económica y continuaron siendo asignadas a la misma comunidad por el algoritmo de Leiden. Incluso la invasión rusa de Ucrania iniciada en 2022 no modificó su pertenencia conjunta a una misma comunidad en ese año. Sin embargo, sería necesario examinar en un estudio posterior cómo evolucionaron los vínculos productivos entre Rusia y Ucrania después de 2022, con datos más actuales que pudieran integrarse en una nueva versión de la \gls{icio}--\gls{ocde}. No obstante, a nivel sectorial, destaca que la extracción de minerales metálicos de Ucrania (B07) se integró de manera estable al bloque productivo encabezado por Alemania a partir de 2015.

\begin{figure}[H]
	\centering
	\includegraphics[width=.9\linewidth]{images/communities_membership/RUS.png}
	\caption{Cantidad de sectores económicos (vértices) por año en la comunidad RUS (Rusia).}
	\label{RUS_timeline}
\end{figure}

Por otra parte, entre 1996 y 2010 se identificó una comunidad formada por Grecia y Chipre, a la que se incorporaron Rumania en 1999 y Bulgaria en 2002. A partir de 2011, estas cuatro economías pasaron a formar parte del bloque productivo encabezado por Turquía. Sin embargo, entre 2013 y 2020, la composición de esta comunidad experimentó frecuentes variaciones. Finalmente, desde 2021, Turquía, Grecia, Chipre y Bulgaria quedaron integrados en la comunidad dirigida por Rusia (véase la figura \ref{ComunidadesPorPais_Rusia}). Esta trayectoria sugiere que, después de la crisis económica de 2008, Grecia perdió gradualmente su liderazgo regional frente a Turquía y que, en los últimos años del periodo estudiado, el bloque encabezado por Rusia amplió sus vínculos económicos con las economías balcánicas y del Mediterráneo oriental. 


\begin{figure}[H]
	\centering
	\includegraphics[width=1\textwidth]{images/ComunidadesPorPais_Rusia.png}
	\caption{Integración de las economías de los países bálticos y antiguas repúblicas soviéticas en comunidades.}
	\label{ComunidadesPorPais_Rusia}
\end{figure}


\begin{figure}[H]
	\centering
	\begin{subfigure}{.83\textwidth}
		\centering
		\includegraphics[width=\textwidth]{images/maps/europe_1995_1.png}
		\caption{1995}
	\end{subfigure}
	\begin{subfigure}{.83\textwidth}
		\centering
		\includegraphics[width=\textwidth]{images/maps/europe_2008_1.png}
		\caption{2008}
	\end{subfigure}
	\begin{subfigure}{.83\textwidth}
		\centering
		\includegraphics[width=\textwidth]{images/maps/europe_2022_1.png}
		\caption{2022}
	\end{subfigure}
	\caption{Mapa de las comunidades identificadas en Europa.$^{*}$}
	\label{Mapa_Comunidades_Europa}
	
	\caption*{\footnotesize $^{*}$ Los sectores económicos se disponen circularmente alrededor del centro del país al que pertenecen y se ordenan en el sentido de las manecillas del reloj. El color de cada sector indica la comunidad a la que pertenece. Las aristas representan únicamente los flujos comerciales cuya magnitud supera el percentil 50 de la distribución regional, y su grosor es proporcional a la intensidad de cada flujo.}
\end{figure}


%\begin{figure}[H]
%	\centering
%	\includegraphics[width=1\textwidth]{sankey_diagrams/europe_deu.png}
%	\caption{Evolución de la pertenencia a comunidades de los sectores económicos de los países de Europa central: Polonia, Alemania, Austria, República Checa, Hungría, Eslovaquia, Eslovenia, Croacia y Rumania.}
%	\label{DiagramaSankey_Europa_DEU}
%\end{figure}


%\begin{figure}[H]
%	\centering
%	\includegraphics[width=1\textwidth]{sankey_diagrams/europe_benelux.png}
%	\caption{Evolución de la pertenencia a comunidades de los sectores económicos de los países de Europa noroccidental: Bélgica, Países Bajos, Luxemburgo, Reino Unido, Irlanda y Malta.}
%	\label{DiagramaSankey_Europa_Benelux}
%\end{figure}


%Aunque en térmnos políticos, jurídicos e institucionales la \gls{efta} se mantuvo independiente de la \gls{cee}, y después, de la \gls{ue}, la integración comercial de los países miembros de la \gls{efta} con la comunidad europea ha sido más profunda que el propio comercio al interior de la asociación.


% EU-EFTA RELATIONS: AN HISTORICAL APPRAISAL
% Evolution of EFTA in Differentiated Integration within Europe 

